% =====================================================================
%  LAB 4 -- 2D-Filter and Image Processing
%  Headings copied word for word from DSP2_Lab4_Exercise.pdf
% =====================================================================
\clearpage
\section{Lab exercise 4: 2D-Filter and Image Processing}

TODO: chapter introduction -- theory about histogram equalisation,
gamma correction and image filtering.

\subsection{Laboratory 4 -- Perform the following procedure of a automatic histogram equalisation:}

\subsubsection{(a) build a histogram of the pixel values}

\begin{figure}[H]
  \includegraphics[width=0.6\textwidth]{lab4-histogram.jpg}
  \caption{histogram of pout.gif pixel values}\label{fig:l4-histogram}
\end{figure}

TODO

\subsubsection{(b) build the cumulative sum of this histogram}

\begin{figure}[H]
  \includegraphics[width=0.6\textwidth]{lab4-cumulative.jpg}
  \caption{cumulative sum of histogram}\label{fig:l4-cumulative}
\end{figure}

TODO

\subsubsection{(c) normalise this cumulative sum to the max. pixel value}

TODO

\subsubsection{(d) use this as a mapping from old-new pixel values (lookup-table)}

TODO

\subsubsection{(e) Test your algorithm on pout.gif show the original picture as well both histograms and the resulting picture.}

\begin{figure}[H]
  \includegraphics[width=0.48\textwidth]{lab4-equalized.jpg}\hspace{4pt}
  \includegraphics[width=0.48\textwidth]{lab4-equalized-comparison.jpg}
  \caption{equalized image and histogram comparison}\label{fig:l4-equalized}
\end{figure}

TODO

\subsubsection{(f) Make a comparison to the histeq function of matlab.}

\begin{figure}[H]
  \includegraphics[width=0.6\textwidth]{lab4-histeq-comparison.jpg}
  \caption{custom equalization vs matlab histeq}\label{fig:l4-histeq}
\end{figure}

TODO

\subsection{Laboratory 4 -- Built your own gamma-correction function by using the template gctemp.m:}

\subsubsection{(a) Test the template with different gammas}

\begin{figure}[H]
  \includegraphics[width=0.6\textwidth]{lab4-gamma-pout.jpg}
  \caption{gamma correction on pout.gif}\label{fig:l4-gamma-pout}
\end{figure}

\begin{figure}[H]
  \includegraphics[width=0.6\textwidth]{lab4-gamma-pout-stretch.jpg}
  \caption{gamma correction with contrast stretch on pout.gif}\label{fig:l4-gamma-pout-stretch}
\end{figure}

TODO

\subsubsection{(b) as well with different pictures}

\begin{figure}[H]
  \includegraphics[width=0.6\textwidth]{lab4-gamma-church.jpg}
  \caption{gamma correction on church2.png}\label{fig:l4-gamma-church}
\end{figure}

\begin{figure}[H]
  \includegraphics[width=0.6\textwidth]{lab4-gamma-church-stretch.jpg}
  \caption{gamma correction with contrast stretch on church2.png}\label{fig:l4-gamma-church-stretch}
\end{figure}

\begin{figure}[H]
  \includegraphics[width=0.6\textwidth]{lab4-gamma-adler.jpg}
  \caption{gamma correction on adler.jpg}\label{fig:l4-gamma-adler}
\end{figure}

\begin{figure}[H]
  \includegraphics[width=0.6\textwidth]{lab4-gamma-adler-stretch.jpg}
  \caption{gamma correction with contrast stretch on adler.jpg}\label{fig:l4-gamma-adler-stretch}
\end{figure}

TODO

\subsection{Laboratory 4 -- Perform different techniques on lowandsp.gif}

\subsubsection{(a) to be finally able to highlight only the membranes and structures of the cells}

\begin{figure}[H]
  \includegraphics[width=0.6\textwidth]{lab4-cells-smoothing.jpg}
  \caption{original, mean filter, and median filter}\label{fig:l4-cells-smooth}
\end{figure}

\begin{figure}[H]
  \includegraphics[width=0.6\textwidth]{lab4-cells-histeq.jpg}
  \caption{histogram equalization applied after filtering}\label{fig:l4-cells-histeq}
\end{figure}

\begin{figure}[H]
  \includegraphics[width=0.48\textwidth]{lab4-cells-edges.jpg}
  \caption{edge detection with laplacian and LoG after equalization}\label{fig:l4-cells-edges}
\end{figure}

TODO

\subsection{Laboratory 4 -- Manipulate the following picture adler.jpg}

\subsubsection{(a) so that the background has only one colour without disturbing the coin}

TODO

\subsubsection{(b) so that the coin has only one colour without disturbing the background}

\begin{figure}[H]
  \includegraphics[width=0.6\textwidth]{lab4-adler-histogram.jpg}
  \caption{adler.jpg original and pixel histogram}\label{fig:l4-adler-hist}
\end{figure}

\begin{figure}[H]
  \includegraphics[width=0.6\textwidth]{lab4-adler-threshold.jpg}
  \caption{adler.jpg thresholding result}\label{fig:l4-adler-thresh}
\end{figure}

TODO

\subsection{Homework 4 (Diary only: Adapt your own histogram equalisation algorithm to perform a good result for the coloured picture church2.png)}

\begin{figure}[H]
  \includegraphics[width=0.6\textwidth]{hw4-hsv-equalization.jpg}
  \caption{HSV-based histogram equalization on church2.png}\label{fig:l4-hw4}
\end{figure}

TODO
